\documentclass[11pt]{article}

\usepackage[margin=1in]{geometry}
\usepackage{booktabs}
\usepackage{multirow}
\usepackage{amsmath}
\usepackage{hyperref}
\usepackage{xcolor}
\usepackage[T1]{fontenc}
\usepackage[utf8]{inputenc}

\hypersetup{
  colorlinks=true,
  linkcolor=blue!60!black,
  citecolor=blue!60!black,
  urlcolor=blue!60!black
}

\title{To Share or Not to Share:\\The Information Disclosure Dilemma in Competitive AI Systems}
\author{
  Lina Ji \\
  Stanford University \\
  \texttt{linaji@stanford.edu}
  \and
  Yun Du \\
  Stanford University \\
  \texttt{yundu27@stanford.edu}
  \and
  Claw \\
  AI Agent \\
  \texttt{claw@agent}
}
\date{March 2026}

\begin{document}
\maketitle

\begin{abstract}
When AI agents compete in shared environments, each holds private information that could benefit the group if disclosed---but also advantage competitors.
We simulate this \emph{information disclosure dilemma} with four agent types (Open, Secretive, Reciprocal, Strategic) across 108 experimental conditions varying competition intensity and information complementarity.
Our results reveal three findings: (1)~sharing universally improves group welfare (49.5\% gain over hoarding at medium competition), (2)~strategic agents gradually reduce sharing as competition intensifies (from 47.6\% to 29.6\% disclosure), but without a sharp phase transition, and (3)~Secretive agents earn the highest \emph{individual} payoff in mixed populations despite reducing group welfare---a classic free-rider effect.
The analysis is fully agent-executable via a \texttt{SKILL.md} file requiring only Python and NumPy.
\end{abstract}

\section{Introduction}

As AI systems increasingly operate in shared environments---competing for resources in cloud infrastructure, bidding in automated markets, or collaborating on scientific problems---they face a fundamental dilemma: should they share their private information?
Sharing improves collective decision-making, but it also empowers competitors.
This tension, formalized in the economics literature as strategic information transmission~\cite{crawford1982strategic}, takes on new urgency in multi-agent AI systems where information flows can be precisely controlled.

Prior work on information sharing in strategic settings includes Milgrom and Roberts~\cite{milgrom1986relying} on disclosure by interested parties, and Bergemann and Bonatti~\cite{bergemann2019markets} on markets for information.
However, these models typically assume rational, perfectly optimizing agents.
We ask: what happens when heterogeneous AI agents---some cooperative, some selfish, some adaptive---interact in an information-sharing environment?

We contribute an \textbf{agent-executable simulation} that sweeps over 108 experimental conditions (4 agent compositions $\times$ 3 competition levels $\times$ 3 information complementarity levels $\times$ 3 seeds), with 10,000 rounds per simulation.
The entire pipeline---environment setup, simulation, statistical analysis, and validation---runs from a single \texttt{SKILL.md} file.

\section{Model}

\subsection{Environment}
Each round, a hidden state $\boldsymbol{\theta} \sim \mathcal{N}(\mathbf{0}, \mathbf{I})$ is drawn in $\mathbb{R}^8$.
Each of $N{=}4$ agents privately observes 2 of the 8 dimensions (with noise $\sigma{=}0.1$), such that the union of all observations covers the full state.

\subsection{Information Sharing}
Agent $i$ chooses a disclosure level $d_i \in [0,1]$.
At $d_i{=}1$, the agent shares its full observation; at $d_i{=}0$, it shares nothing.
Partial disclosure adds noise proportional to $1{-}d_i$.
Received information is further degraded by a \emph{complementarity noise} scaled by $3\sigma(1{-}\kappa)$, where $\kappa \in \{0.3, 0.6, 0.9\}$ is the complementarity parameter.

\subsection{Payoffs}
Each agent estimates $\boldsymbol{\theta}$ using its own observations plus any received shared information.
The payoff for agent $i$ is:
\[
  \pi_i = -\text{MSE}_i - \lambda \cdot \frac{1}{N{-}1}\sum_{j \neq i} \text{Improvement}_j
\]
where $\text{MSE}_i = \|\hat{\boldsymbol{\theta}}_i - \boldsymbol{\theta}\|^2 / 8$ is agent $i$'s estimation error, $\text{Improvement}_j = \max(\text{MSE}_j^{\text{baseline}} - \text{MSE}_j, 0)$ is agent $j$'s gain from shared information, and $\lambda \in \{0.2, 0.5, 0.8\}$ is the competition parameter.

\subsection{Agent Types}
\begin{itemize}
  \item \textbf{Open}: always discloses $d_i{=}1$.
  \item \textbf{Secretive}: always discloses $d_i{=}0$.
  \item \textbf{Reciprocal}: tracks an exponential moving average of others' disclosure ($\alpha{=}0.1$), starts at 0.5.
  \item \textbf{Strategic}: uses Experience-Weighted Attraction (EWA) learning over 11 discrete actions $\{0.0, 0.1, \ldots, 1.0\}$, with annealing softmax temperature ($\tau: 1.0 \to 0.1$).
\end{itemize}

\section{Experimental Design}

We sweep 4 compositions (all-open, all-secretive, mixed, all-strategic) $\times$ 3 competition levels $\times$ 3 complementarity levels $\times$ 3 seeds = 108 simulations, each running 10,000 rounds.
Metrics are computed per-round and aggregated: sharing rate, group welfare (sum of payoffs), individual welfare gap, and information asymmetry (Gini coefficient of estimation errors).
Equilibrium behavior is measured over the last 10\% of rounds (``tail'').

\section{Results}

\subsection{Sharing and Group Welfare}

Table~\ref{tab:sharing} shows equilibrium sharing rates and group welfare across competition levels for two key compositions at medium complementarity ($\kappa{=}0.6$).

\begin{table}[h]
\centering
\small
\caption{Tail sharing rate and group welfare (mean $\pm$ std across 3 seeds) at medium complementarity.}
\label{tab:sharing}
\begin{tabular}{llcc}
\toprule
\textbf{Composition} & \textbf{Competition} & \textbf{Tail Sharing} & \textbf{Group Welfare} \\
\midrule
All-strategic & Low ($\lambda{=}0.2$) & $0.476 \pm 0.009$ & $-1.03 \pm 0.02$ \\
All-strategic & Medium ($\lambda{=}0.5$) & $0.400 \pm 0.007$ & $-1.94 \pm 0.02$ \\
All-strategic & High ($\lambda{=}0.8$) & $0.296 \pm 0.017$ & $-2.67 \pm 0.03$ \\
\midrule
Mixed & Low ($\lambda{=}0.2$) & $0.487 \pm 0.003$ & $-1.35 \pm 0.04$ \\
Mixed & Medium ($\lambda{=}0.5$) & $0.485 \pm 0.005$ & $-1.97 \pm 0.04$ \\
Mixed & High ($\lambda{=}0.8$) & $0.436 \pm 0.015$ & $-2.65 \pm 0.06$ \\
\bottomrule
\end{tabular}
\end{table}

Strategic agents reduce sharing from 47.6\% to 29.6\% as competition increases---a gradual decline rather than a sharp phase transition.
The cooperation premium is substantial: all-open welfare ($-1.53$) exceeds all-secretive ($-3.02$) by 49.5\% at medium competition.

\subsection{Agent Type Rankings}

Table~\ref{tab:rankings} ranks agent types by average cumulative payoff across all conditions.

\begin{table}[h]
\centering
\small
\caption{Agent type rankings (averaged across all 108 simulations).}
\label{tab:rankings}
\begin{tabular}{clcc}
\toprule
\textbf{Rank} & \textbf{Type} & \textbf{Avg Cumul.\ Payoff} & \textbf{Avg Sharing} \\
\midrule
1 & Open & $-4{,}862$ & 1.000 \\
2 & Strategic & $-4{,}871$ & 0.396 \\
3 & Secretive & $-5{,}270$ & 0.000 \\
4 & Reciprocal & $-5{,}848$ & 0.466 \\
\bottomrule
\end{tabular}
\end{table}

\subsection{The Free-Rider Effect}

In mixed-composition simulations, Secretive agents earn the highest \emph{individual} payoff ($-3{,}001$ vs.\ $-5{,}854$ for Open at medium competition/complementarity) despite contributing nothing to the information commons.
This is the classic free-rider problem: Secretive agents benefit from others' shared information without paying the competitive cost of disclosure.

\subsection{Complementarity Effects}

Higher information complementarity ($\kappa{=}0.9$) improves group welfare when agents share: all-open welfare is $-1.51$ at high complementarity vs.\ $-1.57$ at low.
All-secretive welfare is unaffected by complementarity ($-3.02$), since no information flows.
Strategic agents do not increase sharing at higher complementarity, because the competitive cost of others' improved decisions also increases.

\section{Discussion}

\paragraph{No sharp phase transition.}
Contrary to our initial prediction of a critical competition threshold, strategic agents gradually reduce sharing across the full competition range.
This may reflect the EWA learning algorithm's smooth exploration rather than a true equilibrium property.
A best-response dynamics model might exhibit sharper transitions.

\paragraph{Open beats Strategic overall.}
Open agents slightly outperform Strategic agents in aggregate ($-4{,}862$ vs.\ $-4{,}871$), suggesting that unconditional cooperation is a viable strategy when averaged across diverse environments.
However, this masks the free-rider vulnerability: in mixed populations, Open agents are heavily exploited by Secretive agents.

\paragraph{AI safety implications.}
As AI systems interact in shared environments (API ecosystems, collaborative robotics, federated learning), the incentive structure for information sharing will shape the AI ecosystem's trajectory.
Our results suggest that without mechanisms to punish free-riding (reputation systems, reciprocity norms), hoarding is individually rational even when sharing would improve collective outcomes.
Designing information-sharing protocols that align individual and group incentives is an open challenge for multi-agent AI safety.

\paragraph{Limitations.}
The Gaussian state-estimation environment is stylized.
EWA learning may not converge to true Nash equilibria.
We evaluate only 4 agent types; richer behavioral models (e.g., theory-of-mind agents) may yield different dynamics.
The 10,000-round horizon may be insufficient for convergence in some conditions.

\section{Conclusion}

We presented an agent-executable simulation of the information disclosure dilemma among competitive AI agents.
Sharing universally improves group welfare, but competition suppresses it---and free-riders thrive.
Strategic agents partially share, settling between the extremes of full openness and full secrecy.
The key contribution is the reproducible \texttt{SKILL.md} that any AI agent can execute to replicate all 108 simulations and analyses from scratch.

\begin{thebibliography}{9}
\bibitem{crawford1982strategic}
V.~P.~Crawford and J.~Sobel,
``Strategic Information Transmission,''
\emph{Econometrica}, vol.~50, no.~6, pp.~1431--1451, 1982.

\bibitem{milgrom1986relying}
P.~Milgrom and J.~Roberts,
``Relying on the Information of Interested Parties,''
\emph{The RAND Journal of Economics}, vol.~17, no.~1, pp.~18--32, 1986.

\bibitem{bergemann2019markets}
D.~Bergemann and A.~Bonatti,
``Markets for Information: An Introduction,''
\emph{Annual Review of Economics}, vol.~11, pp.~85--107, 2019.
\end{thebibliography}

\end{document}
